% Ad Atticum 2.19 (ad-atticum-02-19)
% Source: https://raw.githubusercontent.com/PerseusDL/canonical-latinLit/master/data/phi0474/phi057/phi0474.phi057.perseus-lat1.xml
% Fetched by scripts/fetch_latin.py

% Dateline (Perseus): Scr. Romae medio m. Quint. a. 695 (59) .

\ciceroLetterOpener{CICERO ATTICO salutem}

\ciceroSection{1} multa me sollicitant et ex rei publicae tanto motu et ex iis periculis quae mihi ipsi intenduntur et sescenta sunt; sed mihi nihil est molestius quam Statium manu missum; nec meum imperium, ac mátto imperium, nón simultatém meam reverteri saltem! nec quid faciam scio neque tantum est in re quantus est sermo. ego autem ne irasci possum quidem iis quos valde amo; tantum doleo ac mirifice quidem. † cetera in magnis rebus†. minae Clodi contentionesque quae mihi proponuntur modice me tangunt; etenim vel subire eas videor mihi summa cum dignitate vel declinare nulla cum molestia posse. dices fortasse: dignitatis ἅλισ tamquam δρυόσ , saluti, si me amas, consule. me miserum! cur non ades? nihil profecto te praeteriret. ego fortasse τυφλώττω et nimium τῷ .

\ciceroSection{2} scito nihil umquam fuisse tam infame, tam turpe, tam peraeque omnibus generibus, ordinibus, aetatibus offensum quam hunc statum qui nunc est, magis me hercule quam vellem non modo quam putarem. populares isti iam etiam modestos homines sibilare docuerunt. Bibulus in caelo est nec qua re scio, sed ita laudatur quasi unus homo nobis cunctando restituit rem. Pompeius, nostri amores, quod mihi summo dolori est, ipse se adflixit. neminem tenent voluntate; ne metu necesse sit iis uti vereor. ego autem neque pugno cum illa causa propter illam amicitiam neque approbo, ne omnia improbem quae antea gessi; utor via. populi sensus maxime theatro et spectaculis perspectus est; nam gladiatoribus qua dominus qua advocati sibilis conscissi; ludis Apollinaribus Diphilus tragoedus in nostrum Pompeium petulanter invectus est; nóstra miseriá tu es magnus— miliens coactus est dicere; Éandem virtutem ístam veniet témpus cum gravitér gemes totius theatri clamore dixit itemque cetera. nam et eius modi sunt ii versus uti in tempus ab inimico Pompei scripti esse videantur: si neque leges neque mores cogunt—, et cetera magno cum fremitu et clamore sunt dicta. Caesar cum venisset mortuo plausu, Curio filius est insecutus. huic ita plausum est ut salva re publica Pompeio plaudi solebat. tulit Caesar graviter. litterae Capuam ad Pompeium volare dicebantur. inimici erant equitibus qui Curioni stantes plauserant, hostes omnibus; Rosciae legi, etiam frumentariae minitabantur. sane res erat perturbata. equidem malueram quod erat susceptum ab illis silentio transiri, sed vereor ne non liceat. non ferunt homines quod videtur esse tamen ferendum; sed est iam una vox omnium magis odio firmata quam praesidio.

\ciceroSection{4} noster autem Publius mihi minitatur, inimicus est. impendet negotium, ad quod tu scilicet advolabis. videor mihi nostrum illum consularem exercitum bonorum omnium, etiam sat bonorum habere firmissimum. Pompeius significat studium erga me non mediocre; idem adfirmat verbum de me illum non esse facturum; in quo non me ille fallit sed ipse fallitur. Cosconio mortuo sum in eius locum invitatus. id erat vocari in locum mortui. nihil me turpius apud homines fuisset neque vero ad istam ipsam ἀσφάλειαν quicquam alienius. sunt enim illi apud bonos invidiosi, ego apud improbos meam retinuissem invidiam, alienam adsumpsissem. Caesar me sibi vult esse legatum.

\ciceroSection{5} honestior declinatio haec periculi; sed ego hoc non repudio. quid ergo est? pugnare malo. nihil tamen certi. iterum dico utinam adesses! sed tamen si erit necesse, arcessemus. quid aliud? quid? hoc opinor. certi sumus perisse omnia; quid enim ἀκκιζόμεθα tam diu? sed haec scripsi properans et me hercule timide. posthac ad te aut, si perfidelem habebo cui dem, scribam plane omnia aut, si obscure scribam, tu tamen intelleges. in iis epistulis me Laelium, te Furium faciam; cetera erunt ἐν αἰνιγμοῖσ . hic Caecilium colimus et observamus diligenter. edicta Bibuli audio ad te missa. iis ardet dolore et ira noster Pompeius.
